% !TEX root = chapter-standalone.tex

\section{Sub-things}
\label{sec:subthings-algebra}

\todotext{@J: review this section}

A recurring situation in mathematics, and in modeling, is the following: we have some structure---say a \SY{semigroup}, \SY{monoid}, or \SY{group}---and we notice that a \emph{part} of it forms a structure of the same kind in its own right.

For example, consider the \SY{monoid}~$\tup{\reals, +, 0}$, which we might use to model the accumulation of some quantity over time.
If the quantity in question can only take whole-number values (counting items in a warehouse, say), then only the elements of the subset~$\wnumbers \setsubseteq \reals$ are relevant to our model.
Conveniently, the subset~$\wnumbers$ is ``self-sufficient'': adding two whole numbers again yields a whole number, and the neutral element~$0$ is a whole number.
So the additive structure of~$\tup{\reals, +, 0}$ restricts to~$\wnumbers$, making the latter a \SY{monoid} in its own right.
We say that it is a \emph{submonoid} of~$\tup{\reals, +, 0}$.

Recognizing such ``sub-things'' is useful in at least two ways.
First, it is a source of new examples: from one structure, we obtain many new ones, one for each sub-thing.
Second, and more importantly for applications, it lets us model \emph{restrictions}: a subsystem of a system, the feasible subset of all conceivable actions, or the transformations that respect a given constraint.
The whole point is that the restricted model then automatically inherits the algebraic properties of the larger one.

\paragraph{Subsemigroups}

The key requirement for a subset of a \SY{semigroup} to be a ``semigroup within the semigroup'' is that we cannot escape the subset by composing its elements.
This property is called \emph{closure}.

\begin{ctdefinition}[Subsemigroup]\label{def:subsemigroup}
    Let~$\sgrpA = \tup{\sgrpAset, \mtimes}$ be a \SY{semigroup}.
    A \maindef{subsemigroup} of~$\sgrpA$ is:

    \constit

    \begin{enumerate}
        \item A subset~$\sgrpBset \setsubseteq \sgrpAset$.
    \end{enumerate}

    \condit

    \begin{enumerate}
        \item The set~$\sgrpBset$ is closed under the composition operation~$\mtimes$ from~$\sgrpA$:
              \begin{equation}\label{eq:subsemigroup-closure}
                  \prfperiod{
                      \sgrpela \setin \sgrpBset \quad \sgrpelb \setin \sgrpBset
                  }{
                      (\sgrpela \mtimes \sgrpelb) \setin \sgrpBset
                  }
              \end{equation}
    \end{enumerate}
\end{ctdefinition}

\begin{remark}
    \label{rem:subsemigroup-is-semigroup}
    The name ``subsemigroup'' is justified by the following observation.
    If~$\sgrpBset \setsubseteq \sgrpAset$ is closed under~$\mtimes$, then the operation~$\mtimes$ restricts to a binary operation
    \begin{equation}
        \mtimes \colon \sgrpBset \cartprod \sgrpBset \sto \sgrpBset,
    \end{equation}
    and this restricted operation is again \SY{associative}: the associative law holds for all triples of elements of~$\sgrpAset$, so in particular it holds for all triples of elements of~$\sgrpBset$.
    Therefore~$\tup{\sgrpBset, \mtimes}$ is itself a \SY{semigroup}.
\end{remark}

\begin{example}
    \label{exa:evennum-subsemigroup}
    Consider the \SY{semigroup}~$\tup{\natnumbers,+}$ introduced in \cref{exa:natnum-semigroup}.
    The subset~$\sgrpBset \setsubset \natnumbers$ of even natural numbers is a \SY{subsemigroup}: the sum of two even numbers is always even.
\end{example}

\begin{example}[A non-example]
    \label{exa:oddnum-not-subsemigroup}
    Staying with~$\tup{\natnumbers,+}$, the subset of \emph{odd} natural numbers is \emph{not} a \SY{subsemigroup}: the sum of two odd numbers is even, so the subset is not closed under composition.

    This illustrates that being a sub-thing is a genuine condition on a subset---most subsets of a \SY{semigroup} are not \SY{subsemigroups}.
\end{example}

\begin{example}
    Consider the \SY{semigroup}~$\tup{\posReals, \cdot}$ of positive real numbers with multiplication.
    The interval~$\makeset{ \ela \setin \posReals \mid \ela \leq 1 }$ is a \SY{subsemigroup}: the product of two numbers that are at most~$1$ is again at most~$1$.

    As an applied interpretation, think of the elements of this \SY{subsemigroup} as \emph{efficiencies} or \emph{attenuation factors} of components in series (each component transmits at most what it receives).
    The closure property expresses a physically evident fact: chaining lossy components can only ever yield a lossy system.
\end{example}

\begin{example}
    Recall the \SY{semigroup}~$\sgrpBset = \makeset{ \mapa^n \mid n \setin \natnumbers }$ generated by the plant development map~$\mapa$ of \cref{exa:plant-trafo-semigroup}.
    The singleton subset~$\makeset{ \mapa^4 }$ is a \SY{subsemigroup}: we will see in \cref{sec:generators-relations-algebra} that~$\mapa^4 \mtimes \mapa^4 = \mapa^4$, since~$\mapa^4$ is the map sending every state to the dead state, and applying it twice changes nothing further.

    In modeling terms: ``being dead'' is a self-perpetuating situation, and this fact is captured algebraically by a one-element \SY{subsemigroup}.
\end{example}

\paragraph{Submonoids}

For \SY{monoids}, closure under composition is not the whole story: a \SY{monoid} also has a \SY{neutral element} as a constituent, and we require a sub-thing to contain it.

\begin{ctdefinition}[Submonoid]\label{def:submonoids}
    Let~$\monoidAdefinition$ be a \SY{monoid}.
    A \maindef{submonoid} of~$\monoidA$ is:

    \constit

    \begin{enumerate}
        \item A subset~$\monoidBset \setsubseteq \monoidAset$.
    \end{enumerate}

    \condit

    \begin{enumerate}
        \item The set~$\monoidBset$ is closed under the composition operation~$\mtimes$ from~$\monoidA$:
              \begin{equation}
                  \prfsemi{\monela \setin \monoidBset \quad \monelb \setin \monoidBset}{\monela \mtimes \monelb \setin \monoidBset}
              \end{equation}

        \item The \SY{neutral element}~$\idmon_\monoidA \setin \monoidAset$ is an element of~$\monoidBset$.
    \end{enumerate}
\end{ctdefinition}

\begin{remark}
    Just as in \cref{rem:subsemigroup-is-semigroup}, a \SY{submonoid}~$\monoidBset$ of a \SY{monoid}~$\monoidA$ is itself a \SY{monoid}: composition restricts to~$\monoidBset$, associativity is inherited, and~$\idmon_\monoidA$ acts as the \SY{neutral element} of~$\monoidBset$ because it acts neutrally on \emph{all} elements of~$\monoidAset$, hence in particular on those of~$\monoidBset$.
\end{remark}

\begin{example}
    $\tup{\natnumbers, +, 0}$,~$\tup{\wnumbers, +, 0}$, and~$\tup{\ratnumbers, +, 0}$ are all \SY{submonoids} of~$\tup{\reals, +, 0}$.
\end{example}

\begin{example}
    $\tup{\natnumbers, \cdot, 1}$ is a \SY{submonoid} of~$\tup{\wnumbers, \cdot, 1}$.
\end{example}

\begin{example}
    $\makeset{ \ela \setin \natnumbers \mid \ela \text{ is even } }$ is a \SY{submonoid} of~$\tup{\natnumbers, +, 0}$.
    Note that the even numbers include~$0$, the \SY{neutral element}.
\end{example}

\begin{example}[A non-example]
    \label{exa:zero-not-submonoid}
    Consider the \SY{monoid}~$\tup{\natnumbers, \cdot, 1}$ of natural numbers with multiplication.
    The singleton subset~$\makeset{0}$ is closed under multiplication, since~$0 \cdot 0 = 0$.
    In fact, $\tup{\makeset{0}, \cdot, 0}$ is itself a perfectly good \SY{monoid}, with~$0$ as its \SY{neutral element}.
    Nevertheless, $\makeset{0}$ is \emph{not} a \SY{submonoid} of~$\tup{\natnumbers, \cdot, 1}$, because it does not contain the \SY{neutral element}~$1$.

    This example shows why the second condition in \cref{def:submonoids} is stated the way it is: we do not merely ask that the subset can be equipped with \emph{some} \SY{monoid} structure; we ask that it inherits \emph{the given} structure, neutral element included.
\end{example}

\begin{example}
    \label{exa:sub-alphabet-submonoid}
    Recall from \cref{exa:string-monoid} that for any set~$\setA$, the set~$\listsof{\setA}$ of lists of elements of~$\setA$ forms a \SY{monoid} under concatenation, with the empty list as \SY{neutral element}.
    If~$\setB \setsubseteq \setA$ is any subset, then the set of lists using only elements of~$\setB$ is a \SY{submonoid} of~$\listsof{\setA}$: concatenating two such lists again yields a list with elements only from~$\setB$, and the empty list trivially qualifies.

    As an applied subexample, let~$\setA$ be the set of all keys of a keyboard, so that elements of~$\listsof{\setA}$ are the possible sequences of keystrokes.
    If~$\setB \setsubseteq \setA$ is the set of keys that are enabled in a given input field (say, only digits, for a field expecting a phone number), then the keystroke sequences that the field accepts form a \SY{submonoid} of all keystroke sequences.
\end{example}

\paragraph{Subgroups}

Continuing the pattern, a sub-thing of a \SY{group} should inherit \emph{all} the constituents of the \SY{group} structure: composition, the \SY{neutral element}, and also inverses.

\begin{ctdefinition}[Subgroup]\label{def:subgroup}
    Let~$\grpA = \tup{\grpAset, \gtimes, \idgrp, \ginv}$ be a \SY{group}.
    A \maindef{subgroup} of~$\grpA$ is:

    \constit

    \begin{enumerate}
        \item A subset~$\grpBset \setsubseteq \grpAset$.
    \end{enumerate}

    \condit

    \begin{enumerate}
        \item The set~$\grpBset$ is closed under the composition operation~$\gtimes$ from~$\grpA$:
              \begin{equation}
                  \prfsemi{
                      \grpela \setin \grpBset \quad \grpelb \setin \grpBset
                  }{
                      \grpela \gtimes \grpelb \setin \grpBset
                  }
              \end{equation}

        \item The \SY{neutral element}~$\idgrp \setin \grpAset$ is an element of~$\grpBset$;

        \item The set~$\grpBset$ is closed under the inverse operation~$\ginv$ from~$\grpA$:
              \begin{equation}
                  \prfperiod{
                      \grpela \setin \grpBset
                  }{
                      \ginv(\grpela) \setin \grpBset
                  }
              \end{equation}
    \end{enumerate}
\end{ctdefinition}

\begin{example}
    The whole numbers~$\wnumbers$ form a \SY{subgroup} of~$\tup{\ratnumbers, +, 0}$, which in turn is a \SY{subgroup} of~$\tup{\reals, +, 0}$.
    Similarly, the even whole numbers form a \SY{subgroup} of~$\tup{\wnumbers, +, 0}$: sums of even numbers are even, $0$ is even, and the negative of an even number is even.
\end{example}

\begin{example}[A non-example]
    \label{exa:natnum-not-subgroup}
    The subset~$\natnumbers \setsubseteq \wnumbers$ is closed under addition and contains the \SY{neutral element}~$0$, so it is a \SY{submonoid} of~$\tup{\wnumbers, +, 0}$.
    However, it is \emph{not} a \SY{subgroup}: it is not closed under the inverse operation, since, for instance, $\ginv(1) = -1 \notsetin \natnumbers$.

    In modeling terms: if group elements represent reversible actions (deposits and withdrawals on an account, say), then~$\natnumbers$ corresponds to allowing only deposits---a consistent restriction of the system, but one in which reversibility is lost.
\end{example}

\begin{example}
    Recall the Klein four-group~$\makeset{I, V, H, R}$ of symmetries of a rectangle from \cref{exa:grp-Klein4}.
    The subset~$\makeset{I, R}$, consisting of ``doing nothing'' and ``rotating by 180 degrees'', is a \SY{subgroup}: composing rotations by 180 degrees yields the identity ($R \gtimes R = I$), and each of the two elements is its own inverse.
    Likewise, $\makeset{I, V}$ and~$\makeset{I, H}$ are \SY{subgroups}.

    The subset~$\makeset{I, V, H}$, on the other hand, is not a \SY{subgroup}: $V \gtimes H = R$, which escapes the subset.
\end{example}

\begin{example}
    \label{exa:matrix-subgroups}
    The \SY{matrix groups} of \cref{sec:matrix-grps} form a chain of \SY{subgroups}:
    \begin{equation}
        \mgSOn \setsubseteq \mgOn \setsubseteq \mgGLn ,
    \end{equation}
    and also~$\mgSOn \setsubseteq \mgSLn \setsubseteq \mgGLn$.

    For instance, to see that~$\mgOn$ is a \SY{subgroup} of~$\mgGLn$: the product of two orthogonal matrices is orthogonal, the identity matrix is orthogonal, and the inverse of an orthogonal matrix (namely its transpose) is orthogonal.

    These \SY{subgroups} are important in engineering because they encode restrictions to transformations that preserve additional structure: matrices in~$\mgOn$ preserve lengths and angles, and matrices in~$\mgSOn$ additionally preserve orientation.
    A rigid body in space can undergo exactly the transformations of~$\mathrm{SO}(3, \reals)$ (rotations); allowing mirror images as well corresponds to enlarging the \SY{subgroup} to~$\mathrm{O}(3, \reals)$.
\end{example}

\begin{example}
    Consider the cyclic group of order~4 from \cref{exa:grp-Z4}, with underlying set~$\makeset{0,1,2,3}$ and composition given by addition modulo~4.
    The subset~$\makeset{0, 2}$ is a \SY{subgroup}: $2 + 2 = 0 \pmod 4$, the \SY{neutral element}~$0$ is included, and each element is its own inverse.

    Thinking of~$\makeset{0,1,2,3}$ as the four positions of a rotary switch, the \SY{subgroup}~$\makeset{0,2}$ corresponds to only allowing half-turns of the switch.
\end{example}

\begin{exercise}
    \label{ex:subgroups-of-Z4}
    Find \emph{all} \SY{subgroups} of the cyclic group of order 4 (\cref{exa:grp-Z4}).
\end{exercise}
\begin{solution}
    There are three: the trivial \SY{subgroup}~$\makeset{0}$, the \SY{subgroup}~$\makeset{0,2}$, and the whole group~$\makeset{0,1,2,3}$.
    No other subset containing~$0$ is closed under addition modulo 4: any subset containing~$1$ or~$3$ generates all four elements by repeated addition.
\end{solution}

\paragraph{Relating the three notions}

Every \SY{group} is in particular a \SY{monoid}, and every \SY{monoid} is in particular a \SY{semigroup}.
It is instructive to note how the three notions of sub-thing interact with this hierarchy.
Each notion adds one closure condition on top of the previous one: a \SY{submonoid} is a \SY{subsemigroup} which contains the \SY{neutral element}, and a \SY{subgroup} is a \SY{submonoid} which is furthermore closed under inverses.
\Cref{exa:zero-not-submonoid,exa:natnum-not-subgroup} show that each of these successive conditions is a genuine further restriction: there are \SY{subsemigroups} that are not \SY{submonoids}, and \SY{submonoids} (of a group) that are not \SY{subgroups}.

\begin{remark}
    For \SY{groups}, there is a well-known compact criterion: a \emph{non-empty} subset~$\grpBset \setsubseteq \grpAset$ is a \SY{subgroup} if and only if
    \begin{equation}
        \prfperiod{
            \grpela \setin \grpBset \quad \grpelb \setin \grpBset
        }{
            \grpela \gtimes \ginv(\grpelb) \setin \grpBset
        }
    \end{equation}
    We leave the verification as an exercise below.
    While handy in proofs, this criterion is arguably less transparent than the three conditions of \cref{def:subgroup}, which spell out, constituent by constituent, what is being inherited.
\end{remark}

\begin{exercise}
    \label{ex:subgroup-criterion}
    Prove the criterion stated in the remark above.
\end{exercise}
\begin{solution}
    First suppose~$\grpBset$ is a \SY{subgroup}.
    Given~$\grpela, \grpelb \setin \grpBset$, closure under inverses gives~$\ginv(\grpelb) \setin \grpBset$, and closure under composition then gives~$\grpela \gtimes \ginv(\grpelb) \setin \grpBset$.

    Conversely, suppose the criterion holds and~$\grpBset$ is non-empty, say~$\grpela \setin \grpBset$.
    Taking~$\grpelb = \grpela$ in the criterion yields~$\grpela \gtimes \ginv(\grpela) = \idgrp \setin \grpBset$.
    Given any~$\grpelb \setin \grpBset$, taking the pair~$\tup{\idgrp, \grpelb}$ in the criterion yields~$\idgrp \gtimes \ginv(\grpelb) = \ginv(\grpelb) \setin \grpBset$, so~$\grpBset$ is closed under inverses.
    Finally, given~$\grpela, \grpelb \setin \grpBset$, we know~$\ginv(\grpelb) \setin \grpBset$, and applying the criterion to the pair~$\tup{\grpela, \ginv(\grpelb)}$ yields~$\grpela \gtimes \ginv(\ginv(\grpelb)) = \grpela \gtimes \grpelb \setin \grpBset$, so~$\grpBset$ is closed under composition.
\end{solution}
